\ifdefined\NoetherPaperTwentyTwoSFourBodyOnly
\providecommand{\NoetherPaperTwentyTwoSFourEnd}{}
\else
\documentclass[11pt]{article}
\usepackage[a4paper,margin=1.45cm]{geometry}
\usepackage[T1]{fontenc}
\usepackage[utf8]{inputenc}
\usepackage{lmodern}
\usepackage[french]{babel}
\usepackage{amsmath,amssymb,mathtools}
\usepackage[protrusion=true,expansion=false]{microtype}
\setlength{\parindent}{1.25em}
\setlength{\parskip}{0.10em}
\allowdisplaybreaks
\setlength{\emergencystretch}{12em}
\sloppy\hbadness=10000\vbadness=10000\hfuzz=2pt\vfuzz=10pt\raggedbottom
\makeatletter\renewcommand\footnotesize{\@setfontsize\footnotesize{7.5}{8.7}\selectfont}\makeatother
\providecommand{\Amod}{\mathfrak A}
\providecommand{\Bmod}{\mathfrak B}
\providecommand{\Gmod}{\mathfrak G}
\providecommand{\Mmod}{\mathfrak M}
\providecommand{\Nmod}{\mathfrak N}
\providecommand{\Rmod}{\mathfrak R}
\providecommand{\Smod}{\mathfrak S}
\providecommand{\mideal}{\mathfrak m}
\providecommand{\nideal}{\mathfrak n}
\providecommand{\gideal}{\mathfrak g}
\providecommand{\hideal}{\mathfrak h}
\providecommand{\jideal}{\mathfrak j}
\providecommand{\tideal}{\mathfrak t}
\providecommand{\aideal}{\mathfrak a}
\providecommand{\bideal}{\mathfrak b}
\providecommand{\bb}{\mathfrak b}
\providecommand{\cc}{\mathfrak c}
\providecommand{\zero}[1]{\equiv 0\pmod{#1}}
\providecommand{\Norm}{N}

\begin{document}
\providecommand{\NoetherPaperTwentyTwoSFourEnd}{\end{document}}
\fi

% Unit: Paper 22 section 4. Direct French translation from the German final-audited source slice.
\setcounter{footnote}{11}

\section*{\S 4. Rapport entre les idéaux de polynômes et les modules de formes linéaires.}

Pour concevoir un idéal $\mideal$ de polynômes, que l'on supposera toujours transformé, comme un module $\Mmod_{i-1}$ $(i=1,\ldots,n)$ de formes linéaires, il faut regarder les polynômes individuels $F(x)$ comme des formes linéaires dans les produits de puissances $\xi_\lambda$ de $x_1\ldots x_{i-1}$:
\[
  F(x)=\sum a^{(i)}_\lambda\xi_\lambda,
\]
où, suivant la notation du \S~3, les $a^{(i)}_\lambda$ désignent des polynômes en $x_i\ldots x_n$. Il résulte aussitôt de la définition de l'idéal que $\Mmod_{i-1}$ possède la propriété de module relativement à ces grandeurs entières $a^{(i)},b^{(i)},\ldots$; en outre, comme les infiniment nombreux produits de puissances $\xi$ de $x_1\ldots x_{i-1}$ n'interviennent que \emph{linéairement}, ils jouent, en raison de leur indépendance linéaire, le rôle d'indéterminées pour $\Mmod_{i-1}$. Chaque module $\Mmod_{i-1}$ contient infiniment beaucoup d'indéterminées $\xi$, mais dans chaque forme linéaire particulière il n'intervient qu'un nombre fini d'indéterminées. De même, chaque idéal fondamental $\gideal_{i-1}$ doit être conçu comme un module $\Gmod_{i-1}$ de formes linéaires, où les indéterminées sont de nouveau les produits de puissances de $x_1\ldots x_{i-1}$. Pour $i=1$, il ne s'agit que d'\emph{une} indéterminée $\xi_0$, qui correspond à l'unité.

On peut maintenant, et c'est là-dessus que repose tout ce qui suit, se limiter, malgré ces infiniment nombreuses indéterminées $\xi$, aux modules de formes linéaires en un nombre fini d'indéterminées, d'après le théorème suivant.

\medskip
\noindent\textbf{Théorème VII.} \emph{Le système des classes de $\Gmod_{i-1}$ modulo $\Mmod_{i-1}$ est isomorphe au système des classes de $\Gmod^*_{i-1}$ modulo $\Mmod^*_{i-1}$, où $\Mmod^*_{i-1}$ désigne un module en un nombre fini d'indéterminées et $\Gmod^*_{i-1}$ son module fondamental. Le module $\Mmod_{i-1}$ ne possède, après adjonction de $x_{i+1}\ldots x_n$ à $P(u)$, qu'un nombre fini de diviseurs élémentaires différents de $1$, à savoir les diviseurs élémentaires de $\Mmod^*_{i-1}$ qui sont différents de $1$.}

Les diviseurs élémentaires de $\Mmod_{i-1}$ doivent ici être définis comme dans la note 8). L'isomorphie qui figure dans le Théorème VII repose sur la définition suivante.

\medskip
\noindent\textbf{Définition V.} \emph{Deux systèmes de classes modulo leurs modules respectifs s'appellent isomorphes lorsqu'ils peuvent être mis en correspondance biunivoque de telle sorte qu'à la différence de deux classes corresponde la différence des classes correspondantes, et qu'au produit d'une classe par une grandeur entière corresponde le produit de la classe correspondante par la même grandeur entière.}

La démonstration du Théorème VII s'obtient par induction complète. Pour $i=1$, le Théorème VII est manifestement exact; car $\Mmod_0$ est un module de formes linéaires en une indéterminée $\xi_0$, qui correspond au produit de puissances de dimension zéro des $x$, donc à l'unité; les grandeurs entières sont tous les polynômes en $x_1\ldots x_n$. L'idéal fondamental $\gideal_0$ devient égal à l'idéal-unité formé de tous les polynômes en $x_1\ldots x_n$, et $\Gmod_0$ est donc égal au module $(\xi_0)$, le module fondamental de $\Mmod_0$, de sorte qu'ici $\Gmod^*_0\mid\Mmod^*_0$ coïncide avec $\Gmod_0\mid\Mmod_0$. Enfin, comme $\Mmod_0$ est de rang $1$, il ne peut posséder au plus qu'un diviseur élémentaire différent de $1$.

Nous passons d'abord de $i=1$ à $i=2$, afin de conduire ensuite la conclusion inductive générale au moyen de l'intuition ainsi obtenue sur la structure de $\Gmod_1\mid\Mmod_1$. Pour cela, nous montrons d'abord que, pour $\Gmod_0\mid\Mmod_0$, on peut prendre un système de représentants borné en $x_1$; en raison de la divisibilité de $\gideal_1$ par $\gideal_0$, cela conduira alors aux indéterminées en nombre fini de $\Gmod^*_1$.

Comme idéal transformé, $\mideal$ possède d'après le \S~3 au moins un polynôme $C^{(1)}(x)$ régulier en $x_1$, disons de dimension $k$; $\Mmod_0$ contient donc la forme linéaire $C^{(1)}\xi_0$. En vertu de la régularité de $C^{(1)}(x)$, tout polynôme $H(x)$ admet une représentation:
\begin{equation}
 H(x)\equiv b^{(2)}_0+b^{(2)}_1x_1+\cdots+b^{(2)}_{k-1}x_1^{k-1}\quad (C^{(1)}(x));
\tag{21}
\end{equation}
ainsi, à cause de l'appartenance de $C^{(1)}\xi_0$ à $\Mmod_0$, on peut effectivement choisir pour $\Gmod_0\mid\Mmod_0$ un système de représentants qui n'atteint pas la $k$-ième dimension en $x_1$.

Si maintenant, en particulier, on pose pour les polynômes $G(x)$ de $\gideal_1$ et $F(x)$ de $\mideal$:
\begin{equation}
 \left\{\begin{aligned}
  G(x)&\equiv g^{(2)}_0+g^{(2)}_1x_1+\cdots+g^{(2)}_{k-1}x_1^{k-1} &&(C^{(1)}(x)),\\
  F(x)&\equiv a^{(2)}_0+a^{(2)}_1x_1+\cdots+a^{(2)}_{k-1}x_1^{k-1} &&(C^{(1)}(x)),
 \end{aligned}\right.
\tag{22}
\end{equation}
et si $\Gmod^*_1$ désigne le système de toutes les formes linéaires $g(\xi)=g^{(2)}_0\xi_0+\cdots+g^{(2)}_{k-1}\xi_{k-1}$, tandis que $\Mmod^*_1$ désigne de façon correspondante le système des formes linéaires $a(\xi)=a^{(2)}_0\xi_0+\cdots+a^{(2)}_{k-1}\xi_{k-1}$, il résulte de la propriété d'idéal de $\gideal_1$ et de $\mideal$ que $\Gmod^*_1$ et $\Mmod^*_1$ sont des modules de formes linéaires en $\xi_0\ldots\xi_{k-1}$ relativement aux grandeurs entières $b^{(2)},c^{(2)},\ldots$. De même, l'idéal principal dérivé de $C^{(1)}(x)$ donne un module relativement à ces grandeurs entières:
\[
 \mathfrak C_1=(\zeta_0,\zeta_1,\ldots,\zeta_\nu,\ldots),
\]
où l'on a posé $\zeta_\nu=x_1^\nu C^{(1)}(x)$. Les $\zeta_\nu$ sont linéairement indépendants relativement à ces grandeurs entières $b^{(2)},c^{(2)},\ldots$, aussi bien entre eux que vis-à-vis des $\xi$ en nombre fini. De la divisibilité de $\mathfrak C_1$ par $\Mmod_1$, et par conséquent par $\Gmod_1$, il résulte que $\Gmod^*_1$ est divisible par $\Gmod_1$, respectivement $\Mmod^*_1$ par $\Mmod_1$. En tenant compte de (22), on obtient donc:
\begin{equation}
 \Gmod_1=(\Gmod^*_1,\mathfrak C_1),\qquad
 \Mmod_1=(\Mmod^*_1,\mathfrak C_1).
\tag{23}
\end{equation}

De (23) et de l'indépendance linéaire des $\xi$ et des $\zeta$, le Théorème VII pour $i=2$ résulte comme suit. À cause de la divisibilité de $\Gmod_1$ par $\Mmod_1$, on a d'abord $\Gmod_1\mid\Mmod_1=\Gmod^*_1\mid\Mmod_1$. Pour démontrer l'isomorphie avec $\Gmod^*_1\mid\Mmod^*_1$, supposons qu'un système de représentants de $\Gmod^*_1\mid\Mmod^*_1$ soit donné par certaines formes linéaires $g^*(\xi)$ de $\Gmod^*_1$. En raison de l'indépendance linéaire des $\xi$ et des $\zeta$, on déduit alors de $g^*(\xi)\equiv 0(\Mmod_1)$ que $g^*(\xi)\equiv 0(\Mmod^*_1)$; les $g^*(\xi)$ sont donc aussi incongrus modulo $\Mmod_1$, forment un système de représentants de $\Gmod_1\mid\Mmod_1$, et la correspondance de $\Gmod^*_1\mid\Mmod_1$ vers $\Gmod^*_1\mid\Mmod^*_1$ transmise par ce système de représentants est isomorphe d'après la Définition V.

Tout à fait de même, on obtient: de $b^{(2)}g(\xi)\equiv0(\Mmod_1)$, $b^{(2)}\ne0$, on déduit que $b^{(2)}g(\xi)\equiv0(\Mmod^*_1)$, $b^{(2)}\ne0$. Comme cela est satisfait pour tous les $g(\xi)$ de $\Gmod^*_1$ et seulement pour ceux-ci -- car $\Gmod^*_1$ consiste, d'après (23), en la totalité des polynômes de l'idéal fondamental $\gideal_1$ qui sont en même temps des formes linéaires en $\xi$ --, $\Gmod^*_1$ est ainsi caractérisé comme module fondamental de $\Mmod^*_1$.

Si l'on adjoint enfin $x_3\ldots x_n$ à $P(u)$, on obtient:
\begin{equation}
\left\{\begin{aligned}
 \Gmod^*_1&=(\eta_1,\ldots,\eta_\rho),&
 \Mmod^*_1&=(e_1\eta_1,\ldots,e_\rho\eta_\rho),\\
 \Gmod_1&=(\eta_\rho,\ldots,\eta_1,\zeta_0,\ldots,\zeta_\nu,\ldots),&
 \Mmod_1&=(e_\rho\eta_\rho,\ldots,e_1\eta_1,\zeta_0,\ldots,\zeta_\nu,\ldots),
\end{aligned}\right.
\tag{24}
\end{equation}
et donc:
\[
 (e_\rho)=\Mmod_1/\Gmod_1=\Mmod_1/(\Mmod_1,\eta_\rho),\qquad
 (e_{\rho-1})=(\Mmod_1,\eta_\rho)/\Gmod_1\quad\hbox{etc.}
\]
Ainsi, en tenant compte de la note 8), $e_\rho,e_{\rho-1},\ldots,e_1,1,\ldots,1,\ldots$ sont reconnus comme les diviseurs élémentaires de $\Mmod_1$. Le Théorème VII est donc démontré en toutes ses parties pour $i=2$.

Il faut remarquer que (22) et (23) n'utilisent que le fait que $C^{(1)}(x)$ est régulier en $x_1$. Ces formules, et les considérations qui s'y rattachent et conduisent au Théorème VII, restent donc valables lorsque, à la place de (12), on prend pour base la transformation spéciale
\begin{equation}
 y_1=x_1;\quad y_2=u_{21}x_1+z_2;\quad\ldots;\quad y_n=u_{n1}x_1+z_n
\tag{25}
\end{equation}
qui, composée avec
\begin{equation}
\left\{\begin{array}{l}
 z=U_1(x)\quad\hbox{ou}\\[2mm]
 z_2=x_2;\quad z_3=u_{32}x_2+x_3;\quad\ldots;\quad z_n=u_{n2}x_2+\cdots+u_{n,n-1}x_{n-1}+x_n
\end{array}\right.
\tag{26}
\end{equation}
redonne (12). Si $\widetilde{\mideal}_{(u)}$ passe, en vertu de (25), en $\widetilde{\mideal}$, et si $\widetilde{\gideal}_1,\widetilde{\Gmod}_1,\ldots$ ont pour $\widetilde{\mideal}$ la même signification que $\gideal_1,\Gmod_1,\ldots$ ont pour $\mideal$, alors on a donc:
\begin{equation}
 \widetilde{\Gmod}_1=(\widetilde{\Gmod}^{*}_1,\widetilde{\mathfrak C}_1),\qquad
 \widetilde{\Mmod}_1=(\widetilde{\Mmod}^{*}_1,\widetilde{\mathfrak C}_1).
\tag{27}
\end{equation}

Ici (27) naît de (23) simplement par l'inversion de (26). Car cela est manifestement vérifié pour $\mideal$ et $C^{(1)}(x)$, donc aussi pour $\mathfrak C_1$ et $\Mmod_1$. Cela vaut aussi pour $\gideal_1$, puisque
\[
 b^{(2)}(x)G(x)\equiv0(\mideal),\quad b^{(2)}\ne0,
 \qquad
 \widetilde b^{(2)}(z)\widetilde G(x,z)\equiv0(\widetilde\mideal),\quad \widetilde b^{(2)}\ne0
\]
s'entraînent réciproquement en vertu de $z=U_1(x)$. Ainsi $\widetilde\gideal_1=\gideal_1$ en vertu de (26); donc aussi $\widetilde\Gmod_1=\Gmod_1$, et, par suite, d'après la définition de $\Gmod_1^*$ et de $\Mmod_1^*$ donnée par (22), il en va de même pour $\widetilde\Gmod_1^*$ et $\widetilde\Mmod_1^*$, et donc pour toute la représentation (27).

Pour la conclusion inductive générale, supposons satisfaites les hypothèses:
\begin{equation}
\left\{\begin{aligned}
 \Gmod_{i-1}&=(\Gmod_{i-1}^*,\mathfrak C_{i-1}),&
 \Mmod_{i-1}&=(\Mmod_{i-1}^*,\mathfrak C_{i-1}),\\
 \mathfrak C_{i-1}&=(\zeta_0,\zeta_1,\ldots,\zeta_r,\ldots),
\end{aligned}\right.
\tag{28}
\end{equation}
où $\Gmod^*_{i-1}$ et $\Mmod^*_{i-1}$ désignent des modules -- relativement aux grandeurs entières $b^{(i)},c^{(i)},\ldots$ -- de formes linéaires en un nombre fini d'indéterminées $\xi$, certains produits de puissances de $x_1,\ldots,x_{i-1}$; et où les $\zeta$ sont linéairement indépendants, aussi bien entre eux que vis-à-vis des $\xi$, relativement à ces grandeurs entières. Une représentation correspondant à (28), ainsi que l'indépendance linéaire des $\xi$ et des $\zeta$, doit également valoir lorsqu'à la place de (12) on prend pour base la transformation spéciale
\begin{equation}
\begin{aligned}
 y_1&=x_1;\quad \ldots;\quad y_{i-1}=u_{i-1,1}x_1+\cdots+x_{i-1};\\
 y_i&=u_{i,1}x_1+\cdots+u_{i,i-1}x_{i-1}+z_i;\quad\ldots;\quad
 y_n=u_{n,1}x_1+\cdots+u_{n,i-1}x_{i-1}+z_n
\end{aligned}
\tag{29}
\end{equation}
qui peut être composée pour donner (12) au moyen de
\[
 z=U_{i-1}(x)\quad\hbox{ou}\quad
 z_i=x_i;\quad z_{i+1}=u_{i+1,i}x_i+x_{i+1};\quad\ldots;\quad
 z_n=u_{n,i}x_i+\cdots+u_{n,n-1}x_{n-1}+x_n,
\]
et cette représentation spéciale doit sortir de (28) simplement par l'inversion de $z=U_{i-1}(x)$.

Des hypothèses (28) et de l'indépendance linéaire des $\xi$ et des $\zeta$ résulte l'isomorphie de $\Gmod_{i-1}\mid\Mmod_{i-1}$ avec $\Gmod^*_{i-1}\mid\Mmod^*_{i-1}$, par correspondance avec le même système de représentants, au moyen des mêmes considérations exactement que celles qui se rattachaient ci-dessus à (23); de même, il en résulte que $\Gmod^*_{i-1}$ devient module fondamental de $\Mmod^*_{i-1}$, et que $\Mmod_{i-1}$ ne possède qu'un nombre fini de diviseurs élémentaires différents de $1$, à savoir les diviseurs élémentaires de $\Mmod^*_{i-1}$ qui sont différents de $1$.

Pour effectuer la conclusion inductive, il faut d'abord montrer de nouveau que, pour $\Gmod^*_{i-1}\mid\Mmod^*_{i-1}$, donc aussi pour $\Gmod_{i-1}\mid\Mmod_{i-1}$, on peut prendre un système de représentants borné en $x_i$. Cela résulte de la représentation du quotient démontrée en (7), Théorème III:
\begin{equation}
 \Gmod^*_{i-1}=\Mmod^*_{i-1}/(C^{(i)}),
\tag{30}
\end{equation}
où $C^{(i)}$ désigne un déterminant quelconque, non identiquement nul, de degré $\rho$ de $\Mmod^*_{i-1}$, $\rho$ étant le rang de $\Mmod^*_{i-1}$. De la partie des hypothèses relative à la transformation spéciale (29), il résulte que $C^{(i)}$ peut toujours être supposé régulier relativement à $x_i$. En effet, le module $\widetilde\Mmod^*_{i-1}$ correspondant à $\Mmod^*_{i-1}$ a le même rang; si donc $\widetilde C^{(i)}$ désigne un déterminant quelconque, non identiquement nul, de degré $\rho$ de $\widetilde\Mmod^*_{i-1}$ qui, par $z=U_{i-1}(x)$, passe en un $C^{(i)}$ régulier, alors, par hypothèse, $C^{(i)}$ devient un déterminant de degré $\rho$ de $\Mmod^*_{i-1}$.

L'existence de $C^{(i)}$ implique, comme en (8), avec une numérotation convenable des $\xi$, l'appartenance à $\Mmod^*_{i-1}$ de $\rho$ formes linéaires de la forme spéciale
\[
 L_\lambda(\xi)=C^{(i)}\xi_\lambda+c^{(i)}_{\lambda,1}\xi_{\rho+1}+\cdots+c^{(i)}_{\lambda,s-\rho}\xi_s
 \qquad(\lambda=1\ldots\rho).
\]
Toute forme linéaire en $\xi$ peut alors être ramenée à la forme:
\begin{equation}
 H(\xi)=a^{(i)}_1\xi_1+\cdots+a^{(i)}_\rho\xi_\rho+b^{(i)}_1\xi_{\rho+1}+\cdots+b^{(i)}_{s-\rho}\xi_s\quad (L_1(\xi),\ldots,L_\rho(\xi)),
\tag{31}
\end{equation}
où $a^{(i)}_1\ldots a^{(i)}_\rho$ sont bornés en $x_i$ et n'atteignent pas le degré $k$ de $C^{(i)}$. Cette représentation est unique; car de
\[
 a^{(i)}_1\xi_1+\cdots+a^{(i)}_\rho\xi_\rho+b^{(i)}_1\xi_{\rho+1}+\cdots+b^{(i)}_{s-\rho}\xi_s\equiv0\quad (L_1(\xi),\ldots,L_\rho(\xi))
\]
on déduit, par comparaison des coefficients de $\xi_1\ldots\xi_\rho$ et en tenant compte des degrés en $x_i$, l'annulation identique de tous les $a^{(i)}$ et $b^{(i)}$.

Soit maintenant en particulier
\[
 H(\xi)\equiv0(\Gmod^*_{i-1});\qquad\hbox{donc}\qquad C^{(i)}H(\xi)\equiv0(\Mmod^*_{i-1})\quad\hbox{d'après (30).}
\]
On obtient alors, comme dans la démonstration du Théorème III:
\[
 C^{(i)}H(\xi)-\sum_{\tau=1}^{s-\rho}\{C^{(i)}b^{(i)}_\tau-\sum_\lambda a^{(i)}_\lambda c^{(i)}_{\lambda,\tau}\}\xi_{\rho+\tau}
 \equiv0(\Mmod^*_{i-1}),
\]
et donc, puisque, comme dans le Théorème III, les coefficients de $\xi_{\rho+1}\ldots\xi_s$ doivent s'annuler:
\[
 C^{(i)}b^{(i)}_\tau=\sum_\lambda a^{(i)}_\lambda c^{(i)}_{\lambda,\tau}\qquad(\tau=1\ldots s-\rho),
\]
d'où il résulte que les $b^{(i)}_\tau$ sont eux aussi bornés et n'atteignent pas le degré maximal des $c^{(i)}_{\lambda,\tau}$. L'existence d'un système de représentants borné en $x_i$ pour $\Gmod^*_{i-1}\mid\Mmod^*_{i-1}$, c'est-à-dire pour $\Gmod_{i-1}\mid\Mmod_{i-1}$, qui n'atteint pas un certain degré fixe $r$, est ainsi démontrée.

Désignons maintenant par $\vartheta_1\ldots\vartheta_t$ les produits de puissances en nombre fini
\[
 x_i^\alpha\xi_\lambda\quad(\alpha=0,1,\ldots,k-1;\ \lambda=1,2,\ldots,\rho),\qquad
 x_i^\beta\xi_{\rho+\tau}\quad(\beta=0,1,\ldots,r-1;\ \tau=1,2,\ldots,s-\rho)
\]
et rangeons les expressions dénombrablement infinies
\[
 x_i^\gamma L_\lambda\quad(\gamma=0,1,\ldots\text{ in inf.};\ \lambda=1,2,\ldots,\rho),
 \qquad
 x_i^\gamma\zeta_\nu\quad(\gamma,\nu=0,1,\ldots\text{ in inf.})
\]
en une suite simplement infinie $\omega_0,\omega_1,\ldots,\omega_\mu,\ldots$; alors les $\vartheta$ et les $\omega$ sont linéairement indépendants, aussi bien séparément entre eux que les uns des autres, relativement aux grandeurs entières $b^{(i+1)},c^{(i+1)},\ldots$. Car de
\[
 \sum c^{(i+1)}_{\alpha\lambda} x_i^\alpha\xi_\lambda+
 \sum c^{(i+1)}_{\beta\tau} x_i^\beta\xi_{\rho+\tau}+
 \sum d^{(i+1)}_{\gamma\lambda} x_i^\gamma L_\lambda(\xi)+
 \sum e^{(i+1)}_{\gamma\nu}x_i^\gamma\zeta_\nu=0,
\]
chaque somme ne portant que sur un nombre fini de termes, il résulte, à cause de l'indépendance linéaire supposée des $\xi$ et des $\zeta$, ainsi que des $\xi$ entre eux relativement aux grandeurs entières $b^{(i)}$, que tous les $e^{(i+1)}_{\gamma\nu}$ s'annulent. De l'unicité de la représentation (31), on déduit ensuite l'annulation des $c^{(i+1)}_{\alpha\lambda}$ et des $c^{(i+1)}_{\beta\tau}$, puis enfin, en raison de l'indépendance linéaire des $L_\lambda(\xi)$, l'annulation des $d^{(i+1)}_{\gamma\lambda}$.

Si l'on conçoit maintenant le module $(\Gmod_{i-1},L_1(\xi),\ldots,L_\rho(\xi))$ comme un module $\mathfrak C_i$ relativement aux grandeurs entières $b^{(i+1)}$, alors $\mathfrak C_i$ devient divisible par $\Mmod_i$, à cause de la divisibilité de $\Gmod_{i-1},L_1(\xi)\ldots L_\rho(\xi)$ par $\Mmod_{i-1}$, et l'on obtient
\[
 \mathfrak C_i=(\omega_0,\omega_1,\ldots,\omega_\mu,\ldots).
\]
Pour tout $H(x)\equiv0(\gideal_{i-1})$, on a donc, d'après (28), (31) et ce qui vient d'être démontré:
\[
 H(x)\equiv b^{(i+1)}_1\vartheta_1+\cdots+b^{(i+1)}_t\vartheta_t\quad(\mathfrak C_i)
\]
comme équation de modules relativement aux grandeurs entières $b^{(i+1)},c^{(i+1)}$. Or $\gideal_i$ est divisible par $\gideal_{i-1}$, et $\mideal$ divisible par $\gideal_i$; si donc, en particulier, on pose pour tout $G(x)$ de $\mathfrak C_i$ et tout $F(x)$ de $\Mmod_i$:
\[
 G(x)\equiv g^{(i+1)}_1\vartheta_1+\cdots+g^{(i+1)}_t\vartheta_t\quad(\mathfrak C_i),\qquad
 F(x)\equiv a^{(i+1)}_1\vartheta_1+\cdots+a^{(i+1)}_t\vartheta_t\quad(\mathfrak C_i),
\]
et si le module formé de tous les $g(\vartheta)$ est désigné par $\Gmod_i^*$, celui formé de tous les $a(\vartheta)$ par $\Mmod_i^*$, les mêmes considérations exactement que celles qui conduisaient à (23) donnent:
\begin{equation}
 \Gmod_i=(\Gmod_i^*,\mathfrak C_i),\qquad
 \Mmod_i=(\Mmod_i^*,\mathfrak C_i),\qquad
 \mathfrak C_i=(\omega_0,\omega_1,\ldots,\omega_\mu,\ldots).
\tag{32}
\end{equation}
Dans la dérivation de (32), seule la régularité de $C^{(i)}$ en $x_i$ a de nouveau été utilisée; toute la considération reste donc valable si l'on prend pour base la transformation spéciale (29) avec un indice augmenté d'une unité. Alors les mêmes considérations que celles qui se rattachent à (27) montrent que cette représentation spéciale sort de (32) simplement par l'inversion de $z=U_i(x)$.

Ainsi toutes les hypothèses (28) etc. de la conclusion inductive générale sont démontrées pour l'indice suivant; et puisque, comme cela a été montré là, le Théorème VII découle immédiatement de ces hypothèses, le Théorème VII est donc démontré en général. En même temps, il est montré que l'arbitraire de $(\Gmod^*_{i-1},\Mmod^*_{i-1})$ donné par l'arbitraire des $C^{(i)}$ est à nouveau supprimé par l'isomorphie du système des classes de $\Gmod^*_{i-1}$ modulo $\Mmod^*_{i-1}$.

En particulier, si le corps $P$ pris au \S~3 comme domaine de coefficients est de caractéristique nulle -- c'est-à-dire si le corps premier contenu dans $P$ et dérivé de l'unité est du type des nombres rationnels --, les indéterminées $u_{\mu\nu}$ peuvent être spécialisées en des grandeurs $\bar u_{\mu\nu}$ de $P$ de telle sorte que, pour $i=1\ldots n$, chaque $C^{(i)}$ reste régulier en $x_i$. Les considérations ci-dessus montrent que le Théorème VII subsiste aussi sous cette spécialisation.\footnote{Hentzelt prend, au lieu d'un $P$ quelconque, seulement le corps de tous les nombres complexes et traite principalement ce dernier cas, tandis qu'il ne prend des indéterminées pour base que dans la théorie de l'élimination proprement dite. Hentzelt montre ensuite, essentiellement par les conclusions données ici, qu'il suffit, pour former la forme résultante, d'aller dans chaque $\Mmod_{i-1}$ seulement jusqu'à un certain degré fini dépassant une borne fixe dans les $x$. Il manque cependant l'isomorphie, donnée dans le Théorème VII, de $\Gmod_{i-1}\mid\Mmod_{i-1}$ avec $\Gmod^*_{i-1}\mid\Mmod^*_{i-1}$, qui expose la raison de cette indépendance à l'égard des nombres de degrés. (E. N.)} Le module fondamental et les diviseurs élémentaires n'ont cependant pas nécessairement à provenir du cas général par la spécialisation $u_{\mu\nu}=\bar u_{\mu\nu}$.\footnote{Pour $\mideal=(x^2,ux+y)$, on a $\gideal_0=(1)$, $\gideal_1=(1)$. Si l'on choisit $C^{(1)}=x^2$, alors $\Gmod_1^*=(1,x)$, $\Mmod_1^*=(ux+y,xy)$; ce dernier $\Mmod_1^*$ possède les diviseurs élémentaires $y^2$ et $1$, et l'on peut choisir $C^{(2)}=y^2$. Pour $u=0$, $C^{(1)}$ et $C^{(2)}$ restent réguliers respectivement en $x$ et en $y$; mais $\Mmod_1^*=(y,xy)$ possède les diviseurs élémentaires $y,y$. Donc $\Gmod_1\mid\Mmod_1$ devient encore isomorphe au système des classes modulo un module fini, mais n'est pas isomorphe à $\Gmod_1\mid\Mmod_1$. Si, en revanche, on avait choisi $C^{(1)}=ux+y$ et spécialisé de façon que $C^{(1)}$ fût resté régulier, alors l'isomorphie aurait elle aussi été conservée. (E. N.)}

\clearpage
\NoetherPaperTwentyTwoSFourEnd
